\begin{abstract}
비원형 체인링은 크랭크 각도마다 체인을 당기는 팔의 길이를 바꾸어, 같은 속도로 달려도 페달이 도는 빠르기를 각도마다 다르게 만든다. 본 연구는 다리를 고관절만 힘을 내는 두 마디 막대로, 근육을 빨라질수록 힘이 줄어드는 직선 관계로 모델링하고, 한 바퀴 중 고관절 부담이 가장 큰 순간을 가장 작게 만드는 체인링 모양을 찾았다. 체인링 운동방정식을 정의와 법칙에서 유도하고, 최적화 문제가 볼록 문제임을 이용해 모양을 푸리에 급수로 일반화하여 정확히 풀었다. 지름비 2, 최소 곡률반지름 25 mm 조건에서 타원형 체인링은 원형보다 가장 큰 부담을 60 rpm에서 1.63\%, 100 rpm에서 2.99\% 줄였고, 위쪽 체인 장력은 한 바퀴 내내 양수였다. 무릎을 힘을 내지 않는 관절로 두면 정강이 힘은 정강이 방향으로만 전달되므로, 페달 힘을 돌리는 성분과 크랭크 방향 성분으로 나누어 선행연구의 효과 지수와 비교하였다.
\end{abstract}

\noindent\textbf{주제어:} 비원형 체인링, 모멘트팔, 고관절 부담, 힘-속도 관계, 볼록 최적화
